% =====================================================================
\section{Proposed Scheme}\label{sec:scheme}
% =====================================================================

\subsection{System Architecture}
Fig.~\ref{fig:arch} shows the architecture of GarbageGo. It consists of
nine modules. The planned deployment is a modular monolith backed by
PostgreSQL/PostGIS for spatial queries.

\begin{figure}[htbp]
\centerline{\figureorplaceholder{architecture.pdf}{\columnwidth}}
\caption{System architecture of GarbageGo. Solid arrows show data flow;
dashed arrows show tool calls issued by the LLM agent.}
\label{fig:arch}
\end{figure}

\begin{itemize}
  \item \modnum{1} \emph{User interface} accepts natural-language requests
        and an authorized location, and displays maps and explanations.
  \item \modnum{2} \emph{LLM agent} identifies intent, calls tools, and
        explains their results. It never generates numbers that a tool
        has not returned.
  \item \modnum{3} \emph{External data and rules} include the vehicle
        feed~\cite{dataset83558}, collection points and schedules~\cite{dataset84004},
        the pedestrian network, and holiday announcements.
  \item \modnum{4} \emph{Data acquisition and validation} normalizes and
        checks timestamps, coordinates, and plates, and records both
        observation time and retrieval time.
  \item \modnum{5} \emph{GIS engine} finds candidate points and estimates
        network walking time with a conservative, user-adjustable speed.
  \item \modnum{6} \emph{ETA engine} estimates the truck's remaining time to
        each point from schedules, live positions, or history, together
        with data age and error information.
  \item \modnum{7} \emph{Domain knowledge} supplies collection rules, time
        slots, and holiday changes. It does not replace ETA or routing.
  \item \modnum{8} \emph{Recommendation engine} removes ineligible or
        clearly missed points, scores the rest by slack, risk, and
        walking effort, and searches for alternatives.
  \item \modnum{9} \emph{Result presentation} shows the recommended point,
        alternatives, data timestamps, and limitations, and states
        clearly when data is insufficient.
\end{itemize}

\subsection{Time-Feasibility Formulation}
For a query issued at time $T$ and each candidate point $s$, all terms
are in minutes:
\begin{equation}
\mathrm{slack}(s) = \mathrm{ETA}(s) - W(s) - B(s)
\label{eq:slack}
\end{equation}
where $\mathrm{ETA}(s)$ is the estimated time from $T$ until the truck
reaches $s$, $W(s)$ is the network walking time from the user to $s$, and
$B(s)$ is a safety buffer. A positive slack means that $s$ is feasible
under the current assumptions; it does not guarantee that the user will
catch the truck.

The buffer reflects three sources of uncertainty:
\begin{equation}
B(s) = b_0 + \alpha\,\Delta t(s) + \beta\,n(s)
\label{eq:buffer}
\end{equation}
where $\Delta t(s)$ is the age of the latest GPS fix of the assigned
truck, $n(s)$ is the number of stops remaining before $s$, and $b_0$
covers walking-speed variation. The term in $n(s)$ follows the observation
that delays propagate to later stops~\cite{tcepb2022report,chen2019travel}.
The coefficients $b_0$, $\alpha$, and $\beta$ will be calibrated on
collected data. Missing the truck usually costs more than walking a bit
farther, so the buffer is deliberately conservative; tolerance also
varies, as 21.2\% of residents in a 2012 survey accepted an arrival error
of at most $\pm5$ minutes~\cite{tc2012survey}. Slack is then mapped
to three levels, \emph{feasible}, \emph{at risk}, and \emph{not
recommended}, instead of a probability, because the system does not yet
have enough data to calibrate probabilities.

When live data is stale or missing, the ETA engine falls back to the
scheduled time and labels the result as \emph{scheduled}, never as a
real-time estimate. If neither a live position nor a schedule can be
matched to $s$, the point is marked as uncertain and is not recommended.

\subsection{Algorithm}
Algorithm~\ref{alg:gg} summarizes the recommendation procedure.

\begin{figure}[htbp]
\noindent\rule{\columnwidth}{0.6pt}\\[-2pt]
\refstepcounter{algcount}\label{alg:gg}%
\textbf{Algorithm~\thealgcount:} Time-feasibility recommendation\\[-6pt]
\noindent\rule{\columnwidth}{0.4pt}
\small
\begin{algorithmic}[1]
\REQUIRE request $Q$, authorized location $L$, query time $T$
\ENSURE ranked points, alternatives, timestamps, and limitations
\STATE parse $Q$; \textbf{if} required inputs are missing \textbf{return} a request for them
\STATE load rules, points, and schedules; \textbf{if} unavailable \textbf{return} data limitation
\STATE $S \leftarrow$ candidate points near $L$; \textbf{if} $S=\emptyset$ widen range or time slot
\STATE fetch live vehicle positions with observation times
\FOR{each $s \in S$}
  \IF{$s$ is ineligible or the truck has clearly passed}
    \STATE remove $s$; \textbf{continue}
  \ENDIF
  \STATE $W \leftarrow$ walking time from $L$ to $s$
  \STATE $E \leftarrow$ ETA of the assigned truck to $s$
  \IF{live data is stale or missing}
    \STATE $E \leftarrow$ scheduled time, labeled \emph{scheduled}
  \ENDIF
  \IF{$W$ or $E$ cannot be determined}
    \STATE mark $s$ uncertain with reason; \textbf{continue}
  \ENDIF
  \STATE $B \leftarrow$ buffer from data age and remaining stops, (\ref{eq:buffer})
  \STATE $\mathrm{slack}(s) \leftarrow E - W - B$; assign feasibility level
\ENDFOR
\STATE rank feasible points by slack, risk, and walking effort
\IF{no feasible point}
  \STATE search other points or the next slot, or report that none exists
\ENDIF
\RETURN recommendation, alternatives, data timestamps, and limitations
\end{algorithmic}
\vspace{-4pt}
\noindent\rule{\columnwidth}{0.4pt}
\end{figure}

\subsection{Role of the LLM Agent}
Following the parser-plus-algorithm split of~\cite{yuan2025llmap}, the
agent converts a request such as ``Where should I take my trash now? I
walk slowly'' into tool parameters, calls the GIS, ETA, and
recommendation engines, and writes an explanation in Traditional Chinese.
Every position, ETA, distance, and walking time in the answer is traceable
to a tool result with its data timestamp. If a tool fails or data is
insufficient, the agent reports the limitation instead of filling in a
value. Because the agent cannot assume the truck will wait, it also never
advises users to run after a truck.
